\documentclass[10pt]{amsart}
\usepackage[a4paper,margin=23mm]{geometry}
\usepackage{amsmath,amssymb,amsthm,mathtools,hyperref,graphicx,xspace,letltxmacro,xpatch}
\usepackage[english]{babel}
\newtheorem{theo}{Theorem}[section]
\newtheorem{lemm}[theo]{Lemma}\newtheorem{prop}[theo]{Proposition}\newtheorem{coro}[theo]{Corollary}\newtheorem{conj}[theo]{Conjecture}
\theoremstyle{definition}\newtheorem{defi}[theo]{Definition}\newtheorem{rema}[theo]{Remark}\newtheorem{exam}[theo]{Example}\newtheorem{exem}[theo]{Example}\newtheorem{nota}[theo]{Notation}
\newcommand{\ie}{i.e.}\newcommand{\eg}{e.g.}\newcommand{\resp}{resp.}\newcommand{\skpt}{\par\smallskip\noindent}
\emergencystretch=3em\pagestyle{plain}
%~Mouliné par MaN_auto v.0.25.0 2022-04-20 09:29:01

\newcommand{\R}{\mathbb{R}}
\newcommand{\N}{\mathbb{N}}

\newcommand{\clE}{\mathcal{E}}

%%%%%%%%%%%%%%%%%%%%%%%%%%%%%%%%%%%%%%%%%%%%%%%%%



\graphicspath{{./figures/}}

\newcommand*{\mk}{\mkern -1mu}
\newcommand*{\Mk}{\mkern -2mu}
\newcommand*{\mK}{\mkern 1mu}
\newcommand*{\MK}{\mkern 2mu}

\hypersetup{urlcolor=purple, linkcolor=blue, citecolor=red}

\let\oldtilde\tilde
\renewcommand*{\tilde}[1]{\mathchoice{\widetilde{#1}}{\widetilde{#1}}{\oldtilde{#1}}{\oldtilde{#1}}}
\let\oldhat\hat
\renewcommand*{\hat}[1]{\mathchoice{\widehat{#1}}{\widehat{#1}}{\oldhat{#1}}{\oldhat{#1}}}
%%%%%%%%%%%%%%%%%%%%%%%%%%%%%%%%%%%%%%%%%%%%%%%%

\title{Global generalized solutions of the reflected one-dimensional Vlasov–Poisson system}
\author{Giorgi, Pierre-Antoine, Nouri, Anne}
\date{doi:10.5802/ahl.134; CC BY 4.0}
\begin{document}
\maketitle
\noindent\textit{Editorial note.} Complete original mathematical body and publisher bibliography retained. Native assumptions and source-local support preserved. Standard cited prerequisites stay cited; no mathematical repairs or extra conditions.\par\medskip
\noindent\textit{Allowed standard prerequisites (editorial).}\begin{itemize}
\item Standard functional analysis, measure theory and the author-cited classical theorems. Standard prior results stay cited; complete source-local proofs and definitions retained.
\end{itemize}\medskip

\section{Introduction}\label{introduction}

We consider the Vlasov--Poisson system
\begin{align}
&\partial_t f+ v \partial_x f+ E\partial_v f= 0, \quad t>0, \quad (x,v) \in\,]0,1[\times \,\R, \label{VP-f}
\\
&\partial_t g+ v \partial_x g- E\partial_v g= 0, \quad t>0, \quad\: (x,v)\in\,]0,1[\times \,\R, \label{VP-g}
\\
&\epsilon \partial_x E= \int_{\R } f\, dv - \int_{\R } g\, dv, \quad t>0,\quad x \in]0,1[, \label{E}
\\
&f(0,x,v)= f^0(x,v), \quad (x,v) \in \,]0,1[\times \,\R, \label{init-f}
\\
&g(0,x,v)= g^0(x,v), \quad\, (x,v) \in \,]0,1[\times \,\R. \label{init-g}
\end{align}
It is used to describe the dynamics of particles in a collisionless, electrostatic and non-relativistic plasma composed of ions and electrons. $f$ and $g$ respectively denote the ionic and electronic distribution functions. The electric field $E$ is created by the ions and electrons themselves and derives from a potential $\phi$ satisfying the Poisson law $\epsilon \partial _{xx}^2\Phi = \int (f-g)\, dv$, i.e~\eqref{E}. The parameter $\epsilon >0$ is equal to the square of the ratio between the Debye and the characteristic observation lengths. The Debye length is a physical length below which charge separation occurs. In many physical situations, $\epsilon $ is small. The distribution functions $f$ and $g$ satisfy direct reflection boundary conditions at the boundary $x \in \{0,1\}$, and $E$ is given and constant at $x=0$,
\begin{align}
f(t,x,v) &= f(t,x,-v),\quad t>0, \quad x\in \{0,1\},\quad v\in \R, \label{boundaries-f}
\\
g(t,x,v) &= g(t,x,-v),\quad t>0, \quad x\in \{0,1\},\quad v\in \R, \label{boundaries-g}
\\
E(t,0) &= E_0,\quad t>0.\label{boundaries-E}
\end{align}
If the spatial domain in~\eqref{VP-f}--\eqref{boundaries-E} were $\R $ instead of $]0,1[$, classical characteristics for~\eqref{VP-f} would be defined by $X^\prime = V, V^\prime = E(t,X)$. In the frame of this paper, bounces may occur at the boundary $\{0,1\}$ of the spatial domain. When a characteristics intersects $x=0$ (resp. $x=1$) with a zero velocity at a time $s>0$, it can be seen that it stays in the domain without any discontinuity. And so, it is still considered as a classical characteristics. Bounces at the boundary of the domain occur when a characteristics intersects the boundary with a non zero velocity. Generalized characteristics involving possible bounces~\cite{Guo95,HS08} and continuous solutions to~\eqref{VP-f} and~\eqref{boundaries-f} are defined as follows.

\begin{defi}\label{generalized_characteristics}
Let $E\in C([0,+\infty [; C^1([0,1]))$. The generalized backward characteristics $(X,V)$ from $(t,x,v) \in\, ]0,+\infty [
\times]0,1[\,\times \R$ related to~\eqref{VP-f} and~\eqref{boundaries-f} is defined as the union of the classical characteristics which connect $(t,x,v)$ to $(t_1,x_1,v_1)$, $(t_1,x_1,-v_1)$ to $(t_2,x_2,v_2), \ldots, (t_n,x_n,-v_n)$ to $(t_{n+1},x_{n+1},v_{n+1}), \ldots$, where $x_n \in
\{0,1\}$, $|v_n|>0$ and $t_n>t_{n+1} \geq 0$.
\end{defi}

This gives a set $P \subset \N ^*$ counting the number of bounces, and a sequence of bouncing times $(t_n)_{n\,\in\,P}$ such that,
\begin{align*}
X\left(t_n;t,x,v\right) = 0 \quad &\text{and}\quad V\left(t_n^+;t,x,v\right) = - V\left(t_n^-;t,x,v\right)> 0,
\\
\intertext{or}
X(t_n;t,x,v) = 1\quad &\text{and}\quad V\left(t_n^+;t,x,v\right) = - V\left(t_n^-;t,x,v\right) < 0,\quad n\in P.
\end{align*}

\begin{defi}\label{df-f}
Let $I$ be an interval of non negative times starting at zero. Let $E\in C(I; C^1([0,1]))$. A continuous solution $f$ to~\eqref{VP-f}, \eqref{init-f} and~\eqref{boundaries-f} on $I$ is a function $f\in C(I\times [0,1] \times \R)$ such that
\begin{equation}\label{continuite-f}
f(t,x,v)= f^0\big(X(0; t,x,v),V(0; t,x,v)\big),\quad t\in I,\quad (x,v)\in [0,1]\times \R,
\end{equation}
where $(X(\cdot ; t,x,v),V(\cdot ; t,x,v))$ is the generalized characteristics from $(t,x,v)$ as in Definition~\ref{generalized_characteristics}.
\end{defi}


The main results of this paper are the following.


\begin{theo}\label{main-result-existence}
Let $\epsilon >0$. Let $f^0, g^0 \in C([0,1] \times \R)$ be non negative even functions with respect to the $v$ variable, with finite kinetic energy, and such that for any $R>0$,
\begin{equation}\label{hyp1-f0}
\begin{split}
\sup_{(x,\,w)\,\in\,[0,\,1]\,\times\,\R ;\,|w-v|\,<\,R}\; &f^0(x,w) \in L^1\left(\R_v\right), \\\sup_{(x,\,w)\,\in\,[0,\,1]\,\times\,\R ; \,|w-v|\,<\,R}\; &g^0(x,w) \in L^1\left(\R_v\right).
\end{split}
\end{equation}
Let $E_0\in \R \setminus \{ 0\} $ such that
\begin{equation}\label{hyp-E0-E1different-de-zero}
E_0 \neq \epsilon ^{-1} \int \left(g^0-f^0\right)(x,v)dxdv.
\end{equation}
There exists a solution
\[
(f,g,E) \in \big(C\big([0,+\infty [\times [0,1] \times \R\big)\big)^2 \times C\big([0,+\infty [; C^1([0,1])\big)
\]
of~\eqref{VP-f}--\eqref{boundaries-E} in the sense of Definition~\ref{df-f}. Moreover, $f$ and $g$ are non negative.
\end{theo}


\begin{theo}\label{main-result-unicite}
Let $f^0, g^0$ be non negative Lipschitz functions satisfying~\eqref{hyp1-f0}, even with respect to the $v$ variable, and such that for some $c_0>0$ and $V_0>0$,
\[
f^0(x,v)= g^0(x,v)= 0,\quad x\in [0,1],\quad |v| \leq c_0 \quad \text{or}\quad |v| \geq V_0.
\]
Let $E_0\in \R \setminus \{ 0\} $ satisfying~\eqref{hyp-E0-E1different-de-zero}. 
There is a time $T_0>0$ such that the continuous solution to~\eqref{VP-f}--\eqref{boundaries-E} is unique on $[0, T_0] $.
\end{theo}

The Vlasov--Poisson system has been studied for long. In order to present the results of this paper with respect to the previous works on the Vlasov--Poisson system, let us give a list of previous works on the subject. 
\begin{itemize}
\item On classical solutions in $(x,v)\in \R ^3\times \R ^3$
\begin{itemize}
\item[-] C.~Bardos and P.~Degond~\cite{BD85} proved global in time existence and uniqueness of classical $C^1$ solutions to the Cauchy problem related to the Vlasov--Poisson system for small initial data.
\item[-] K.~Pfaffelmoser~\cite{Pfa92} proved existence of classical solutions for general initial data in $C^1(\R ^6)$ with compact support. Refinements were done in~\cite{Sch91} and~\cite{Rei97}.
\end{itemize}
\end{itemize}
The proofs for these results are based on an analysis of the characteristics associated to the system.
\vspace{4pt}
\begin{itemize}
\item On weak solutions in $\R ^3\times \R ^3$
\begin{itemize}
\item[-] A.~Arsenev~\cite{Ars75} proved the global existence of $L^1\cap L^\infty $ weak solutions to the Vlasov--Poisson system.
\item[-] Using velocity averages, R.~J.~DiPerna and P.-L.~Lions~\cite{DL88} proved the global existence of renormalized weak solutions.
\item[-] R.~Robert~\cite{Rob97} established uniqueness of weak solutions with compact support.
\item[-] Together with the study of propagation of moments $\int |v|^m f(t,x,v) \, dx dv$ with $m>3$, P.-L.~Lions and B.~Perthame~\cite{LP91} proved existence and uniqueness of an $L^1 \cap L^{\infty}$ weak solution. 
\item[-] C.~Pallard~\cite{Pal12} proved an analogous result for $m>2$.
\item[-] Using optimal transport, G.~Loeper~\cite{Loe06} proved the uniqueness of weak solutions with bounded mass density.
\item[-] E.~Miot~\cite{Mio16} proved uniqueness of weak solutions with the $L^p$ norms of the mass density growing at most linearly with respect to $p$.
\item[-] T.~Holding and E.~Miot~\cite{HM18} proved uniqueness of weak solutions with mass density in an Orlicz space, as a consequence of a quantitative stability result for the Wasserstein distance of two weak solutions.
\end{itemize}
\vspace{4pt}
\item On weak solutions in $[0,+\infty [\times \R $
\begin{itemize}
\item[-] It is known from a counterexample by Y.~Guo~\cite{Guo95} that there is in general no $C^1$ solution to the Vlasov--Maxwell system with direct reflection boundary conditions on a half-line. This counterexample can be adapted to the Vlasov--Poisson system with direct reflection boundary conditions in a bounded domain.
\item[-] For specific frames like species with the same sign of charge or neutral plasmas, Y.~Guo~\cite{Guo94} proved the existence of global in time classical solutions.
\item[-] H.~J.~Hwang and J.~Schaeffer~\cite{HS08} proved uniqueness of weak solutions to the one-species Vlasov--Poisson system with direct reflection for the distribution function and given constant electric field pointing inward of the domain at the boundary. They used an approach with characteristics.
\end{itemize}
\vspace{4pt}
\item On weak solutions in $\Omega \times \R ^3$, where $\Omega $ is a bounded subset of $\R ^3$
\begin{itemize}
\item[-] Existence of weak solutions to the Vlasov--Poisson system in a bounded domain and given indata was proven by R.~Alexandre~\cite{Ale93} and N.~Ben Abdallah~\cite{Ben94}.
\item[-] Existence and stability of weak solutions to the initial boundary value problem was proven by J.~Weckler~\cite{Wec95}.
\item[-] Existence of weak solutions to the corresponding stationary problem was proven by F.~Poupaud~\cite{Pou90}.
\item[-] For the Vlasov equation with a given force field, S.~Mischler~\cite{Mis99} considered other types of boundary conditions such as specular reflection, and proved existence and uniqueness of weak solutions.
\end{itemize}
\end{itemize}


Theorem~\ref{main-result-existence} of this paper proves global in time existence of continuous solutions to the two-species Vlasov--Poisson system in a slab, with specular reflection boundary conditions. It takes into account previous results~\cite{Guo95} about weak solutions in $[0,+\infty [
\times \hspace*{0.02in}\R $, stating that $C^1$ solutions can not be expected. It is new compared to the previous results on weak solutions in $[0,+\infty [\times \hspace*{0.02in}\R $ or $\Omega \times \R ^3$ where $\Omega $ is a bounded subset of $\R ^3$, because of the regularity of its solutions. Indeed the solutions of the Vlasov-Poisson system in Theorem~\ref{main-result-existence} are continuous, i.e. with a regularity between the $C^1$ regularity of strong solutions (in whole space) and the $L^p$ regularity of weak solutions. 


Theorem~\ref{main-result-unicite} of this paper proves local in time uniqueness of the solution found in Theorem~\ref{main-result-existence}. For a fixed electric field, it would concern the previously quoted result~\cite{HS08} by H.~J.~Hwang and J.~Schaeffer. However, the assumption in~\cite{HS08} of an electric field pointing inward the domain at the boundary only allows one bounce at the boundary. In the frame of this paper with two-species of opposite electric charges - which is the case with ions and electrons in most plasmas - if the electric field points inward the domain for a species, it points outward the domain for the other species. It is the reason why the proofs in this paper could not be built on similar arguments as in~\cite{HS08}.


The main tool in the proofs of Theorems~\ref{main-result-existence}-\ref{main-result-unicite} is the use of generalized characteristics introduced in~\cite{Guo95}. Although they enlighten the solutions by following the trajectories of the particles, the main difficulties in their use are to discard the phenomenon of infinitely many bounces accumulating at some boundary point, which is performed in Propositions~\ref{cas-simple}-\ref{nb-fini-rebonds-cas-simple}, and to get continuity of the distribution functions $f$ and $g$ from~\eqref{continuite-f}. Indeed, the presence of boundaries may induce discontinuities at initial time. This problem is overcome by assuming the initial distribution functions even with respect to the $v$-variable.


Here are two remarks on the assumptions on $E_0$ in Theorems~\ref{main-result-existence}-\ref{main-result-unicite}.
\begin{rema}
Due to the mass conservation
\begin{align*}
\int _0^1\int _\R f(t,x,v)dxdv &= \int _0^1\int _\R f^0(x,v)dxdv,\\
\int _0^1\int _\R g(t,x,v)dxdv &= \int _0^1\int _\R g^0(x,v)dxdv, t>0,
\end{align*}
an integration of~\eqref{E} with respect to the space variable implies that
\begin{align*}
&E(t,1)= E_0+\frac{1}{\epsilon }\int _0^1\int _\R \left(f^0-g^0\right)(x,v)dxdv,\quad t>0,
\end{align*}
hence that $E(\cdot,1)$ is a constant function. Assumption~\eqref{hyp-E0-E1different-de-zero} is satisfied if $E_0\neq 0$ and electroneutrality holds at $t=0$. It is done for $E(\cdot,1)$ to be different from zero.
\end{rema}
\begin{rema}
As will be seen in Theorem~\ref{electroneutralite}, $t \mapsto E(t,0)$ is assumed to be a constant function in order to ensure the conservation of total energy. However, Theorems~\ref{main-result-existence}-\ref{main-result-unicite} would also hold for non constant and non vanishing continuous $t\rightarrow E(t,0)$, such that $t\rightarrow E(t,0)+\frac{1}{\epsilon }\int _0^1\int _\R (f^0-g^0)(x,v)dxdv$ does not vanish.
\end{rema}

The paper organizes as follows. In Section~\ref{generalized-characteristics}, generalized characteristics are considered, taking into account possible bounces, proving that infinitely many bounces can not occur, and studying their regularity. Theorem~\ref{main-result-existence} (resp. Theorem~\ref{main-result-unicite}) is proven in Section~\ref{pf-main-result-existence} (resp. Section~\ref{pf-main-result-unicite}). The quasineutrality equation when $\epsilon
\to 0$ is proven in Section~\ref{th-electroneutralite}.

\section{Generalized characteristics}\label{generalized-characteristics}

In this section, $T>0$ and $E\in C([0,T] ; C^1([0,1]))$ are given. We consider the Cauchy problem for the Vlasov equation,
\begin{align}
\partial_t f+ v \partial_x f+ E\partial_v f &= 0, \mkern90mu t\in [0,T],\quad (x,v) \in\;]0,1[\times \,\R, \label{VP-f-2}
\\
f(0,x,v) &= f^0(x,v), \mkern133mu (x,v) \in\;]0,1[\times \,\R, \label{init-f-2}
\\
f(t,x,v) &= f(t,x,-v), \quad t\in [0,T], \quad (x,v) \in \{ 0,1\} \times \,\R. \label{boundaries-f-2}
\end{align}
As recalled above, existence and uniqueness of weak solutions to the problem have been proven by S.~Mischler in~\cite{Mis99}, using a variety of test functions. Our appr\-oach differs from his by considering continuous solutions and using generalized characteristics. Excluding the case where infinitely many bounces would accumulate at a boundary point, we prove in Propositions~\ref{cas-simple} and~\ref{nb-fini-rebonds-cas-simple} that the backwards in time generalized characteristics from any $(t,x,v)\in [0,T]
\times [0,1] \times \R$,
\[
s \mapsto \big(X(s;t,x,v), V(s;t,x,v)\big),
\]
has a finite number of bounces at the boundary, hence reaches time zero. Example~\ref{ex-discontinuity-V(0)} exhibits a case where the map $(t,x,v)\mapsto V(0;t,x,v)$ is discontinuous.\linebreak In Proposition~\ref{continuite-valeur-absolueV(0)}, the continuity of the map $(x,v) \mapsto (X(0;t,x,v), \mid V(0;t,x,v)\mid)$ is proven, for any $t>0$.

\begin{lemm}\label{V-bounded}
For $t>0$ and $(x,v)\in [0,1] \times \R$, the second component $V$ of the generalized characteristics from $(t,x,v)$ satisfies
\begin{equation}\label{v_bounded}
\left| V(s;t,x,v) \right| \leq |v| + T \| E \|_{\infty},
\end{equation}
for $s\in [0,t]$ if $P$ is finite (resp. $s\in [t_1,t] \cup _{n\,\in\,P}[t_{n+1},t_n] $ if $P$ is not finite).
\end{lemm}

\begin{proof}[{Proof of Lemma~\ref{V-bounded}}] The case where $P$ is finite is classical. Assume $P=\N^*$. Denote by $t_0=t$ and prove by induction on $n
\in \N$ that
\begin{equation}\label{property-lemma-V-bounded}
\big | |v|-|V(s;t,x,v)| \big | \leq (t-s) \| E \|_{\infty }, \quad s \in \left[t_{n+1},t_n\right].
\end{equation}
For $n=0$ and $s \in [t_1,t]$, the equation $\partial_s V(s;t,x,v) = E(s,X(s;t,x,v))$ yields
\[
\big| |v| - |V(s;t,x,v)| \big| \leq \int_{s}^t \big|E(r,X(r;t,x,v))\big| dr \leq (t-s) \| E \|_{\infty }.
\]
Assuming that $| |v| - |V(s;t,x,v)| | \leq (t-s) \| E \|_{\infty }$ for $s \in [t_n,t_{n-1}]$ and $n \geq 1$,\linebreak it holds
\begin{align*}
\big| |v| - |V(s;t,x,v)| \big| &\leq \big| |v| - \left|V\left(t_n^+;t,x,v\right)\right|\big| + \big| \left|V\left(t_n^-;t,x,v\right)\right| - |V(s;t,x,v)|\big| \\
&\leq \left(t-t_n\right) \| E \|_{\infty } + \int_s^{t_n} |E(r,X(r;t,x,v))| dr \\
&\leq (t-s) \| E \|_{\infty }, \quad s \in \left[t_{n+1},t_n\right].\qedhere
\end{align*}
\end{proof}

Given $(t,x,v) \in [0,T] \times [0,1] \times \R$, either the backwards in time generalized characteristics from $(t,x,v)$ reaches $\{ 0\} \times [0,1]
\times \R $ at $(X(0;t,x,v),V(0;t,x,v))$ without any bounce. This is the easy case where $P = \varnothing$. Or bounces occur at $x=0$ or $x=1$. We first prove that there is a finite number of them. This is strongly linked with the sign of the electric field at the boundaries. The following analysis distinguishes two cases. We first deal with the case of a negative value of $t \mapsto E(t,1)$. The case of a positive value of $t \mapsto E(t,0)$ can be treated analogously.


In the case of a negative value of $E(t,1)$ on $[0,T]$, let $\delta \in]0,1[$ be such that
\begin{equation}\label{E-bande}
E(t,x)<0, \quad (t,x) \in [0,T] \times [1-\delta, 1].
\end{equation}
Let
\begin{equation}\label{df-Delta(v)}
\Delta(v) = \frac{1}{\|E\|_{\infty }} \left (\sqrt{ \left (|v| + T \| E \|_{\infty } \right)^2 + 2 \delta \|E\|_{\infty } } - \left (|v| + T \| E \|_{\infty } \right) \right), \quad v \in \R.
\end{equation}
A bound of the number of possible bounces on the generalized characteristics from \linebreak
$(t,x,v)\in [0,T] \times [0,1] \times \R$ is given in the following proposition.


\begin{prop}\label{cas-simple}
Assume $t \mapsto E(t,1)$, $t \in [0,T]$, constant and negative. 
The number of bounces occurring at $x=1$ along the backwards in time generalized characteristics from $(t,x,v)\in [0,T] \times [0,1] \times \R $ is finite and bounded by $\frac{T}{\Delta(v)}$.
\end{prop}


\begin{proof}[{Proof of Proposition~\ref{cas-simple}}]
Let $(X,V)$ be a backwards in time generalized characteristics from $(t,x,v)$ with at least two bounces at $x=1$, occurring at times $t_1$ and $t_2$, with $t_2<t_1<t$. In order to prove the result, it is sufficient to bound $t_1-t_2$ from below.
It holds that $V(t_1^-;t,x,v)>0$. As $V(\cdot;t,x,v)$ is decreasing when $X(\cdot;t,x,v)$ is in the interval $[1-\delta,1]$, $X$ leaves the interval $[1-\delta,1]$ at a time $s_1\in]t_2,t_1[$, and $X(s;t,x,v) \geq 1-\delta$ for $s \in [s_1,t_1]$. The integration between $s$ and $t_1$ along the classical characteristics
\[
\big(X(\cdot ;t,x,v),V(\cdot ;t,x,v)\big)\big |_{\left[t_2^+,\,t_1^-\right]}
\]
yields
\[
X(s) - 1 + \left(t_1-s\right) V\left(t_1^-\right) \geq -\frac{\|E\|_{\infty }}{2} \left(t_1-s\right)^2, \quad s \in \left[s_1, t_1\right],
\]
so that
\[
\frac{\|E\|_{\infty }}{2} \left(t_1-s_1\right)^2 + \left(t_1 - s_1\right) V\left(t_1^-\right) - \delta \geq 0.
\]
This implies that
\[
t_1-t_2 \geq t_1 - s_1 \geq \frac{\sqrt{V\left(t_1^-\right)^2 + 2 \delta \|E\|_{\infty }}-V\left(t_1^-\right)}{\|E\|_{\infty }}.
\]
And so, by Lemma~\ref{V-bounded},
\[
t_1-t_2 \geq \frac{1}{\|E\|_{\infty }} \left (\sqrt{ \left (|v| + T \|E\|_{\infty } \right)^2 +2\delta \|E\|_{\infty } } - \left (|v| + T \|E\|_{\infty } \right)\right).
\]
The number of bounces at $x=1$ on $[0,T]$ is thus smaller than $\frac{T}{\Delta(v)}$.
\end{proof}

The opposite case where $E(t,1)>0$, $t\in [0,T]$, is more complicated. It corresponds to an electric field pointing outward of the domain $]0,1[$ at the boundary. It cannot be directly expected that the time between two bounces is bounded from below. An infinite number of bounces is a priori not impossible. The distance to the boundary of the $X$ component of the characteristics between two bounces could be arbitrarily small. It is proven in the following proposition that this does not occur. We first prove a preliminary lemma.

\begin{lemm}\label{convergent_series}
For any $(t,x,v)\in [0,T] \times [0,1] \times \R $, the series $\sum_{n\,\in\,P} |V(t_n^+;t,x,v)|$ converges.
\end{lemm}

\begin{proof}[{Proof of Lemma~\ref{convergent_series}}]
Assume $P=\N^*$, and first prove that the number of consecutive bounces between $x=0$ and $x=1$ is finite. Integration along the classical characteristics on $[t_{k+1}, t_k]$, with $(k,k+1)\in P^2$, yields
\[
X\left(t_k\right) - X\left(t_{k+1}\right) = \left(t_k - t_{k+1}\right) V\left(t_{k+1}^+\right) + \int_{t_{k+1}}^{t_k} \left(t_k-r\right) E(r,X(r))dr.
\]
It holds that $|X(t_k) - X(t_{k+1})|=1$. Hence
\begin{align*}
1 &\leq \left(t_k - t_{k+1}\right) \left|V\left(t_{k+1}^+\right)\right| + \frac{\|E\|_{\infty }}{2} \left(t_k - t_{k+1}\right)^2,
\\
\intertext{so that}
t_k - t_{k+1} &\geq \frac{1}{\| E \|_{\infty }} \left (\sqrt{V\left(t_{k+1}^+\right)^2 + 2\| E \|_{\infty }} - \left|V\left(t_{k+1}^+\right)\right| \right).
\end{align*}
It follows from Lemma~\ref{V-bounded} that
\begin{align*}
t_k - t_{k+1} &\geq \frac{1}{\| E \|_{\infty }} \left (\sqrt{(|v|+T\|E\|_{\infty })^2 + 2\|E\|_{\infty }} - (|v|+T\|E\|_{\infty }) \right).
\end{align*}
This implies a finite number of consecutive bounces between $x=0$ and $x=1$ on $[0,T]$. Consequently we can assume $X(t_k)=0$ for $k \geq n_0$ for some $n_0 \in \N^*$, the other case $X(t_k)=1$ for $k \geq n_0$ being analogous. It holds that
\begin{multline*}
2 \sum_{k=1}^n \left|V\left(t_k^+\right)\right|\\
\begin{aligned}
&= 2 \left(\sum_{k=1}^{n_0-1}\left|V\left(t_k^+\right)\right| \right) + V\left(t_{n_0}^+\right) + V\left(t_n^+\right) + \sum_{k=n_0}^{n-1} \underbrace{\left(\left|V\left(t_{k+1}^+\right)\right|+\left|V\left(t_k^+\right)\right|\right)}_{=V\left(t_{k+1}^+\right) - V\left(t_k^-\right)} \\
&= 2 \left(\sum_{k=1}^{n_0-1} \left|V\left(t_k^+\right)\right| \right) + V\left(t_{n_0}^+\right) + V\left(t_n^+\right) - \int_{t_n}^{t_{n_0}} E(s,X(s))ds,\quad n \geq n_0+1.
\end{aligned}
\end{multline*}
By~\eqref{v_bounded} applied to $s= t_n^+$,
\begin{align*}
&\sum_{k=1}^n \left|V\left(t_k^+\right)\right| \leq \sum_{k=1}^{n_0-1} \left|V\left(t_k^+\right)\right| + \frac{V\left(t_{n_0}^+\right)+|v|}{2} + T \| E \|_{\infty },\quad n \geq n_0.
\end{align*}
\end{proof}


\begin{prop}\label{nb-fini-rebonds-cas-simple}
When $E(t,1)$ takes a constant value $E_1>0$ on $[0,T]$, the number of bounces occurring at $x=1$ along the backwards in time generalized characteristics from $(t,x,v)\in [0,T] \times [0,1] \times \R $ is finite.
\end{prop}

\begin{proof}[{Proof of Proposition~\ref{nb-fini-rebonds-cas-simple}}]
It is a proof by contradiction. Assume infinitely many bounces at $(t_n,1)_{n\,\in\,\N}$ along the backwards in time generalized characteristics from $(t,x,v)\in [0,T]
\times [0,1] \times \R $. $(t_n)_{n\,\in\,\N}$ being a decreasing sequence in $[0,T]$, converges to a time $t^* \geq 0$ when $n \to +\infty $. Denote by
\[
V_n = V\left(t_n^-\right) > 0, s_n \in\left]t_{n+1},t_n\right[\quad \text{such that}\; s_n= \min_{s\,\in\,\left[t_{n+1},\,t_n\right]} X(s),\; y_n = 1-X\left(s_n\right).
\]
It holds that $\lim_{n\,\to\,+\infty} s_n = t^*$, $\lim_{n\,\to\,+\infty} V_n = 0$ by Lemma~\ref{convergent_series}, and $\lim_{n\,\to\,+\infty} y_n = 0$. Indeed, if for a subsequence $(y_{n_k})_{k\,\in\,\N}$ of $(y_n)$, $\lim_{k\,\to\,+\infty} y_{n_k} = y^*> 0$, then
\begin{multline*}
\frac{\| E\|_\infty }{2}(t_{n_k}-s_{n_k})^2+\left(t_{n_k}-s_{n_k}\right)-\frac{y^*}{2}\\
\begin{aligned}
&\geq \int _{s_{n_k}}^{t_{n_k}}\left(t_{n_k}-r\right)E(r,X(r))dr+V\left(s_{n_k}\right)\left(t_{n_k}-s_{n_k}\right)-\frac{y^*}{2}\\
&= X\left(t_{n_k}\right)-X\left(s_{n_k}\right)-\frac{y^*}{2}= y_{n_k}-\frac{y^*}{2}\geq 0,
\end{aligned}
\end{multline*}
for $k$ large enough. Consequently,
\[
t_{n_k}-s_{n_k}\geq \frac{\sqrt{1+y^*\|E\|_\infty }-1}{\|E\|_\infty },
\]
for $k$ large enough. This would contradict the infinite number of bounces at $x=1$. 
By the continuity of $E$, there is $\delta_1>0$ such that
\[
E(t,x) \geq \frac{E_1}{2}, \quad (t,x) \in [0,T] \times \left[1-\delta_1,1\right].
\]
It follows from
\[
0 = \int_{t_{n+1}}^{t_n} V(s)ds = \left(t_n - t_{n+1}\right) V_n - \int_{t_{n+1}}^{t_n} \left(r - t_{n+1}\right) E(r,X(r))dr,
\]
that
\begin{equation}\label{asymp_t_n-cas-simple}
t_n - t_{n+1} \leq \frac{4}{E_1}V_n,
\end{equation}
for $n$ large enough. Moreover,
\begin{align*}
y_n &= 1-X\left(s_n\right)= \int_{s_n}^{t_n} V(s)ds \\
&= \left(t_n - s_n\right) V_n - \int_{s_n}^{t_n} \left(r-s_n\right) E\left(r,X(r)\right)dr 
\leq \left(t_n-s_n\right)V_n,
\end{align*}
for $n$ large enough. And so,
\begin{equation}\label{control-yn}
y_n \leq \frac{4}{E_1}V_n^2.
\end{equation}
In a neighborhood of $(t^*,1)$, $E(s,X(s))$ expresses as
\begin{equation}\label{dev-E-cas-simple}
E(s,X(s)) = E_1 + (X(s)-1) \partial_x E(t^*,1)+(X(s)-1) \epsilon (s),
\end{equation}
with
\begin{equation}\label{lim-eps}
\lim_{s\,\to\,t^*} \epsilon (s) = 0.
\end{equation}
Indeed, the function $\epsilon $ introduced in~\eqref{dev-E-cas-simple} satisfies
\[
\epsilon (s)= \frac{1}{X(s)-1} \int_1^{X(s)} \big(\partial_x E(s,y) - \partial_x E(t^*,1)\big)dy,
\]
which tends to zero when $s$ tends to $t^*$, by the uniform continuity of $\partial_x E$ on\linebreak $[0,T] \times [0,1]$. The case $\partial_x E(t^*,1)<0$, i.e. $\partial_x E(t^*,1)= - \alpha^2 $ with $\alpha>0$, is treated here. $\alpha $ is taken as $1$ for the sake of simplicity. By definition of the characteristics $(X,V)$ and~\eqref{control-yn}--\eqref{lim-eps}, $s \mapsto (X(s;t,x,v), V(s;t,x,v))$ satisfies
\begin{align*}
X^{\prime \prime}(s) +X(s) &= E_1+ 1+g(s),\quad s\in \left[t_{n+1},t_n\right],\\
V(s) &= X^\prime (s),\mkern77mu s\in\,\left]t_{n+1},t_n\right[,\\
X(t_n) &= 1,\quad X^\prime \left(t_n^-\right)= V_n,
\\
\intertext{where}
g(s) &= (X(s)-1) \epsilon (s)= o\left(V_n^2\right),\quad s\in \left[t_{n+1},t_n\right].
\end{align*}
Here, $o(V_n^2)$ means that $\lim_{n\,\to\,+\infty}\tfrac{o(V_n^2)}{V_n^2}= 0$. Hence,
\begin{align*}
X(s) &= E_1\big(1-\cos \left(s-t_n\right)\big)+ V_n\sin \left(s-t_n\right)+ 1+o\left(V_n^3\right), \quad s\in \left[t_{n+1},t_n\right],\\
V(s) &= E_1\sin \left(s-t_n\right)+ V_n \cos \left(s-t_n\right)+o\left(V_n^2\right),\mkern96mu s\in\,\left]t_{n+1},t_n\right[,
\end{align*}
or
\begin{align}
&X(s)= \frac{E_1}{2}\left(s-t_n\right)^2+V_n\left(s-t_n\right)+1+o\left(V_n^3\right), \quad s\in \left[t_{n+1},t_n\right],\label{ordre-position}
\\
&V(s)= E_1\left(s-t_n\right)+V_n+o\left(V_n^2\right),\mkern122mu s\in\left]t_{n+1},t_n\right[.\label{ordre-vitesse}
\end{align}
Since $t_{n+1}$ is a solution to $X(t_{n+1}) = 1$, $t_n-t_{n+1}$ satisfies
\begin{gather*}
E_1\left(t_n-t_{n+1}\right)^2-2V_n\left(t_n-t_{n+1}\right)+o\left(V_n^3\right)= 0,
\\
\intertext{i.e.}
E_1\left(t_n-t_{n+1}\right) = V_n \pm \sqrt{V_n^2 +o\left(V_n^3\right)}.
\end{gather*}
By definition of $V_{n+1}$,
\[
V_{n+1} = - V\left(t_{n+1}^+\right) = E_1\left(t_n-t_{n+1}\right)-V_n+o\left(V_n^2\right).
\]
Given the positive sign of $V_n$ and $V_{n+1}$, it results
\[
E_1\left(t_n-t_{n+1}\right) = V_n +\sqrt{V_n^2 +o\left(V_n^3\right)}= 2V_n+o\left(V_n^2\right)\quad \text{and}\quad V_{n+1} = V_n+o\left(V_n^2\right).
\]
Thus, for some $n_0$ large enough,
\[
V_{n+1} \geq V_n \left(1 - V_n\right), \quad n \geq n_0.
\]
Denote by $h : x \mapsto x - x^2$ and
\[
h^p = \underbrace{h \circ \cdots \circ h}_{p \text{ times}},\;p \in \N^*.
\]
Let $n_1\geq n_0$ be such that $V_n\leq
\tfrac{1}{2}$, $n\geq n_1$. It holds that
\begin{equation}\label{inequ-V-n}
V_{n_1+p} \geq h^{p}\left(V_{n_1}\right), \quad p \in \N.
\end{equation}
Moreover, it can easily be proven by induction that
\begin{equation}\label{f^(n)}
h^p(x) \geq \frac{x}{p+1},\quad p \geq 1,\quad x \in \left[0, \frac{1}{2} \right].
\end{equation}
It results from~\eqref{inequ-V-n}--\eqref{f^(n)} that
\[
\sum_{p= n_1+1}^n V_p \geq \sum_{p= 1}^{n-n_1} h^p \left(V_{n_1}\right) \geq V_{n_1} \sum_{p= 1}^{n-n_1} \frac{1}{p},\quad n \geq n_1+1,
\]
which contradicts the statement of Lemma~\ref{convergent_series}. Hence the number of bounces at $x=1$ is finite. The case $\partial_x E(t^*,1)>0$ (resp. $\partial_x E(t^*,1)=0$) is similar and also leads to~\eqref{ordre-position}-\eqref{ordre-vitesse} (see~\cite{Gio19}).
\end{proof}

Proposition~\ref{nb-fini-rebonds-cas-simple} can similarly be extended to the case where $t \mapsto E(t,1)$ is not a constant function.

We now consider the continuity of the function $f$ defined by~\eqref{continuite-f}. Despite the continuity of $f^0$, $f$ may be discontinuous. Actually, issues arise when $X(0;t,x,v)$ is exactly zero (or one). This is illustrated in the following example.


\begin{exam}\label{ex-discontinuity-V(0)}
Let the field $E$ be a positive constant. Let $t\in] 0,T]$ and $x>\frac{Et^2}{2}$.

The map $v \mapsto V(0;t,x,v)$ is discontinuous at $v=\frac{x}{t}+\frac{E}{2}t$ and continuous elsewhere.
\end{exam}

\begin{proof}[{Proof of Example~\ref{ex-discontinuity-V(0)}}]
The backward in time characteristics from $(t,x,v)$ before any bounce at $x=0$ is given by
\[
X(s)= x-v(t-s)+\frac{E}{2}(t-s)^2, \quad V(s)= v-E(t-s), \quad s \in [0,t].
\]
For $v=\frac{x}{t}+\frac{E}{2}t$, there is no bounce on $]0,t]$, $X(0;t,x,v)=0$ and $V(0;t,x,v)\linebreak >0$. For $\tilde{v}<v$, the backward in time characteristics from $(t,x,\tilde{v})$ has no bounce and
\[
V\left(0;t,x,\tilde{v}\right)= \tilde{v}-tE.
\]
For $\tilde{v}>v$, the backward in time characteristics from $(t,x,\tilde{v})$ encounters a bounce at time
\begin{gather*}
t_1=t-\frac{\tilde{v}-\sqrt{\tilde{v}^2-2xE}}{E}>0,\quad \text{and}\quad V\left(0;t,x,\tilde{v}\right) = \tilde{v}-tE-2\sqrt{\tilde{v}^2-2xE}.
\\
\intertext{Thus,}
\lim _{\tilde{v}\,\to\,v^-}V\left(0;t,x,\tilde{v}\right)= v-tE, \quad \lim _{\tilde{v}\,\to\,v^+}V\left(0;t,x,\tilde{v}\right)= -(v-tE).
\end{gather*}
Hence the map $v \mapsto V(0;t,x,v)$ is discontinuous at $v=\tfrac{x}{t}+\tfrac{E}{2}t$. 
\end{proof}


\begin{prop}\label{continuite-valeur-absolueV(0)}
The map $(t,x,v) \mapsto (X(0;t,x,v),|V(0;t,x,v)|)$ is continuous on $[0,T]
\times [0,1] \times \R $.
\end{prop}

\begin{proof}[{Proof of Proposition~\ref{continuite-valeur-absolueV(0)}}]
Let $(\tilde{t},\tilde{x},\tilde{v})$ be given. If the backwards characteristics from $(\tilde{t},\tilde{x},\tilde{v})$ reaches $t= 0$ at $X(0;\tilde{t},\tilde{x},\tilde{v})\in]0,1[$, then analogous arguments as for classical characteristics imply that
\[
\lim _{(t,\,x,\,v)\,\rightarrow\,\left(\tilde{t},\,\tilde{x},\,\tilde{v}\right)}\big(X(0;t,x,v), V(0;t,x,v)\big) = \big(X\left(0;\tilde{t},\tilde{x},\tilde{v}\right), V\left(0;\tilde{t},\tilde{x},\tilde{v}\right)\big).
\]
What remains to be proven is the continuity of $(t,x,v) \mapsto (X(0;t,x,v),|V(0;t,x,v)|)$ at $(\tilde{t},\tilde{x},\tilde{v})$ such that its backwards characteristics reaches $t= 0$ at $X(0;\tilde{t},\tilde{x},\tilde{v})\in \{ 0,1\} $. Assume $X(0;\tilde{t},\tilde{x},\tilde{v})= 0$. Consider $(t,x,v)$ such that the backwards characteristics $(X,V)(\cdot;t,x,v)$ has an earliest bounce at time $t_1(t,x,v)>0$. For $(t,x,v)$ close enough to $(\tilde{t},\tilde{x},\tilde{v})$, there is no bounce of the backwards characteristics from $(\tilde{t},\tilde{x},\tilde{v})$ on the interval $[0,t_1(t,x,v)]$. Define the extended electric field $E^e$ on $[0,T] \times
\R $ by
\begin{align*}
E^e(t,x) &= E(t,x),\quad t\geq0,\hspace*{0.02in}x\in [0,1],\\
E^e(t,x) &= E(t,0),\quad t\geq 0,\hspace*{0.02in}x<0,\\
E^e(t,x) &= E(t,1),\quad t\geq 0,\hspace*{0.02in}x>1,
\end{align*}
and the extended classical characteristics $(X^e,V^e)$ from $(t,x,v)\in [0,T]
\times [0,1] \times \R $~by
\[
\partial _sX^e= V^e,\hspace*{0.02in} X^e(t;t,x,v)= x,\quad \quad \partial _sV^e= E^e(s,X^e),\; V^e(t;t,x,v)= v.
\]
%\hspace*{1.in}\\

First consider the case where $E_0<0$. There is no restriction to consider $\tilde{v}>0$, and $(t,x)$ (resp. $(\tilde{t},\tilde{x})$) in the strip close to $x=0$ where $E<0$. 


Let $\epsilon >0$ be given, and $v\in]
\frac{\tilde{v}}{2},\frac{3\tilde{v}}{2}[$. For $(t,x,v)$ in an appropriate neighborhood of $(\tilde{t},\tilde{x},\tilde{v})$, it holds that
\begin{multline}\label{dm-continuite-valeur-absolueV(0)-1}
\left|X^e(s;t,x,v)-X^e\left(s;\tilde{t},\tilde{x},\tilde{v}\right)\right|+\left|V^e(s;t,x,v)-V^e\left(s;\tilde{t},\tilde{x},\tilde{v}\right)\right|\\
\leq \epsilon, \quad s \in \left[0,\tilde{t}\right].
\end{multline}
It follows from
\begin{align*}
&V^e(t_1;t,x,v)\geq v,\quad X^e(t_1;t,x,v)= 0,\\
&X^e(s;t,x,v)= V^e(t_1;t,x,v)(s-t_1)+\int _s^{t_1}(r-s) E^e\left(r,X^e(r;t,x,v)\right)dr,
\end{align*}
that
\[
X^e(s;t,x,v)\leq -2\epsilon, \qquad s \leq t_1 - \frac{4\epsilon }{\tilde{v}}.
\]
Together with~\eqref{dm-continuite-valeur-absolueV(0)-1} and $X(0,\tilde{t},\tilde{x},\tilde{v})= 0$, this implies that $t_1<\frac{4\epsilon }{\tilde{v}}$. Consequently,
\[
\left|V\left(t_1^-;t,x,v\right)-V(0;t,x,v)\right| + \left|V\left(t_1;\tilde{t},\tilde{x},\tilde{v}\right)-V\left(0;\tilde{t},\tilde{x},\tilde{v}\right)\right| \leq 2 \|E\|_{\infty} t_1 \leq c\epsilon,
\]
and
\begin{multline*}
\left| V(0;t,x,v)+V\left(0;\tilde{t},\tilde{x},\tilde{v}\right)\right| \\
\begin{aligned}
&\leq \left| V(0;t,x,v)-V\left(t_1^-;t,x,v\right) \right| + \left| V\left(0;\tilde{t},\tilde{x},\tilde{v}\right)-V\left(t_1;\tilde{t},\tilde{x},\tilde{v}\right)\right|\\
&\mkern243mu+ \left|V\left(t_1;\tilde{t},\tilde{x},\tilde{v}\right)-V\left(t_1^+;t,x,v\right)\right| \\
&\leq c\epsilon.
\end{aligned}
\end{multline*}
The inequality
\[
\left|X(0;t,x,v)-X\left(0;\tilde{t},\tilde{x},\tilde{v}\right)\right| \leq c\epsilon,
\]
can be proven by bounding $|X(0;t,x,v)-X(0;\tilde{t},\tilde{x},\tilde{v})|$ from above by
\begin{align*}
\big|X(0;t,x,v)-X\left(t_1;t,x,v\right)\big| &+ \left|X\left(0;\tilde{t},\tilde{x},\tilde{v}\right)-X\left(t_1;\tilde{t},\tilde{x},\tilde{v}\right)\right|\\
 &+ \left|X\left(t_1;\tilde{t},\tilde{x},\tilde{v}\right)-X\left(t_1;t,x,v\right)\right|.
\end{align*}
Consider the case where $E_0>0$ and $(t,x,v)$ (resp. $(\tilde{t},\tilde{x},\tilde{v})$) such that the backwards characteristics $(X,V)(\cdot;t,x,v)$ (resp. $(X,V)(\cdot;\tilde{t},\tilde{x},\tilde{v})$) has an earliest bounce at time $t_1(t,x,v)>0$ (resp. has no bounce on $[0,t_1(t,x,v)]$ and is such that $X(0;\tilde{t},\tilde{x},\tilde{v})= 0$). There is no restriction to consider $\tilde{v}>0$, and $(t,x)$ (resp. $(\tilde{t},\tilde{x})$) in the strip close to $x=0$ where $E>\frac{E_0}{2}$. Let $\epsilon >0$ be given. It holds that for $(t,x,v)$ in an appropriate neighborhood of $(\tilde{t},\tilde{x},\tilde{v})$,
\[
\left|X^e(s;t,x,v)-X^e\left(s;\tilde{t},\tilde{x},\tilde{v}\right)\right|+\left|V^e(s;t,x,v)-V^e\left(s;\tilde{t},\tilde{x},\tilde{v}\right)\right|\leq \epsilon, \quad s \in \left[0,\tilde{t}\right].
\]
Consider the extreme case where $V(t_1(t,x,v);t,x,v)= 0$.


Then $0 \leq V(s;\tilde{t},\tilde{x},\tilde{v})\leq \epsilon $, for $s\in [0,t_1(t,x,v)]$. And so,
\[
\frac{E_0}{2} t_1\leq \int_0^{t_1} E\left(r,X\left(r;\tilde{t},\tilde{x},\tilde{v}\right)\right)dr = V\left(t_1;\tilde{t},\tilde{x},\tilde{v}\right)-V\left(0;\tilde{t},\tilde{x},\tilde{v}\right) \leq c\epsilon.
\]
From here the proof is analogous to the case where $E_0<0$.
\end{proof}
Consequently, taking $f^0$ continuous and even w.r.t. the $v$ variable and defining
\[
f(t,x,v)=f^0\big(X(0;t,x,v),V(0;t,x,v)\big),\quad (t,x,v)\in [0,+\infty[\,\times\,[0,1] \times \R,
\]
as in~\eqref{continuite-f}, makes $f$ continuous.



\section{Proof of the existence Theorem~\ref{main-result-existence}}\label{pf-main-result-existence}


In this section, the parameter $\epsilon $ in~\eqref{E} does not play any role and is taken as one for the sake of simplicity. Let $T>0$ be given. Theorem~\ref{main-result-existence} is proven with a fixed point argument for the map $S$ defined on
\[
K = \left\{ a \in C([0,T] \times [0,1]) \, ; \int_0^1 a(t,x)dx = \left\|f^0\right\|_{L^1} - \left\|g^0\right\|_{L^1},\quad t\in [0,T] \right \},
\]
by $S= S_3 \circ S_2 \circ S_1$. Here,
\begin{align*}
S_1(a)(t,x) &= E_0 +\int_0^x a(t,y)dy, \quad (t,x)\in [0,T] \times [0,1],\\
S_2(E) &= (f,g),
\end{align*}
where $f$ (resp. $g$) is the solution to the linear Vlasov equation with force field $E$ (resp. $-E$), initial datum $f^0$ (resp. $g^0$) and direct reflection boundary conditions, and
\[
S_3(f,g) = \int (f-g)(\cdot, \cdot, v)dv.
\]
In a first step we prove that $S$ maps $K$ into $K$. In a second step we prove the compactness of $S$ in $C([0,T] \times [0,1])$. In a third step we prove its continuity. By a Schauder theorem we conclude that there is a fixed point $a$ of $S$. The definition of $a$ finally implies that $(f,g)= S_2\circ S_1(a)$ is a solution to the Cauchy problem~\eqref{VP-f}--\eqref{boundaries-E}.


\begin{proof}[{First step of the proof.}]
Let us prove that the map $S$ is well defined and maps $K$ into $K$. $S_1(a) = E$ is continuous on $[0,T] \times [0,1]$ like $a$, and globally Lipschitz with respect to $x$ since
\begin{equation}
\left| E\left(t,x'\right) - E(t,x) \right| \leq \| a \|_\infty |x' - x|, \quad t \in [0,T],\quad (x,x') \in [0,1]^2.\label{E_lip}
\end{equation}
Moreover, $E(t,0)= E_0$ and $E(t,1)= E_0+ \| f^0 \|_{L^1} - \| g^0 \|_{L^1}$ are constants different from zero by~\eqref{hyp-E0-E1different-de-zero}. The analysis from Section~\ref{generalized-characteristics} and the evenness of $(f^0,g^0) $ with respect to $v$ allow to define $S_2(E)= (f,g)$, where $f$ (resp. $g$) is the solution to the linear Vlasov equation with force field $E$ (resp. $-E$), initial datum $f^0$ (resp. $g^0$) and direct reflection boundary conditions. Recall that
\begin{align*}
f(t,x,v)&= f^0\big(X(0;t,x,v),|V(0;t,x,v)|\big),\\
\big(\text{resp. } g(t,x,v)&= g^0\big(Y(0;t,x,v),|W(0;,t,x,v)|\big)\big),\quad (t,x,v)\in [0,T] \times [0,1] \times \R,
\end{align*}
where $(X,V)$ (resp. $(Y,W)$) are the generalized characteristics associated to $E$ (resp. $-E$), and are such that
\begin{equation}\label{continuity-characteristics}
(t,x,v)\rightarrow \big(X(0;t,x,v),\lvert V(0;t,x,v)\rvert,Y(0;t,x,v),\lvert W(0;t,x,v)\rvert\big)
\end{equation}
is continuous. Consequently, $f$ (resp.~$g$) is continuous, and nonnegative like $f^0$ (resp.~$g^0$). 
$S(a)=\int (f-g)(\cdot,\cdot,v)dv$ belongs to $C([0,T] \times [0,1])$. Indeed, let
\[
R= \left(\left|E_0\right| + \|a \|_\infty \right) T.
\]
By~\eqref{hyp1-f0}, there is $U>0$ such that
\[
\int_{|v|\,>\,U}\;\sup_{x\,\in\,[0,\,1],\,|w-v|\,<\,R} f^0(x,w)dv \quad \text{and} \quad \int_{|v|\,>\,U}\; \sup_{x\,\in\,[0,\,1],\,|w-v|\,<\,R} g^0(x,w)dv
\]
are arbitrarily small. It follows from the continuity of~\eqref{continuity-characteristics} and $(f^0,g^0)$, that the map
\begin{multline*}
(t,x) \mapsto \int _{|v|\,<\,U}(f-g)(t,x,v)dv\\
= \int_{|v|<U} \Big(f^0\big(X(0;t,x,v),\lvert V(0;t,x,v)\rvert\big)-g^0\big(Y(0;t,x,v),\lvert W(0;t,x,v)\rvert\big) \Big) dv
\end{multline*}
is continuous on $[0,T] \times [0,1]$. Finally, the mass conservation of $f$ (resp. $g$) implies that
\[
\int_0^1 \int (f-g)(t,x,v)dvdx= \left\|f^0\right\|_{L^1} - \left\|g^0\right\|_{L^1}, \quad t \in [0,T].
\]
Consequently $S(a)$ belongs to $K$.
\end{proof}


\begin{proof}[{Second step of the proof}]
Let us prove that $S$ is compact in $C([0,T] \times [0,1])$. Let $(a_n)_{n\,\in\,\N}$ with $a_n\in C([0,T] \times [0,1])$ bounded by $M$. Denote by
\[
E_n = S_1\left(a_n\right)\quad \text{and} \quad \left(f_n,g_n\right) = \left(S_2 \circ S_1\right)\left(a_n\right).
\]
By~\eqref{hyp1-f0}, $(S(a_n))_{n\,\in\,\N}$ is bounded in $C([0,T] \times [0,1])$ by
\[
\int \sup_{x\,\in\,[0,\,1],\,|w-v|\,<\,\left(\left|E_0\right|\,+\,M\right)T} f^0(x,w)dv + \int \sup_{x\,\in\,[0,\,1],\,|w-v|\,<\,\left(\left|E_0\right|\,+\,M\right)T} g^0(x,w)dv.
\]
Prove its uniform equicontinuity. Let $\eta>0$ be given. By~\eqref{hyp1-f0}, there is $U>0$ such that,
\begin{align*}
&\int_{|v|\,>\,U} \sup_{x\,\in\,[0,\,1],\,|w-v|\,<\,\left(|E_0|\,+\,M\right)T} f^0(x,w)dv\\
+&\int_{|v|\,>\,U} \sup_{x\,\in\,[0,\,1],\,|w-v|\,<\,\left(|E_0|+M\right)T} g^0(x,w)dv < \frac{\eta }{2}.
\end{align*}
Let $(X_n,V_n)$ (resp. $(Y_n,W_n))$ be the generalized characteristics associated to the field $S_1(a_n)$ (resp. $-S_1(a_n))$. The existence of $h_0>0$ such that
\begin{align*}
&\sup _{(t,\,x)\,\in\,[0,\,T]\,\times\,[0,\,1]} \int_{|v|\,<\,U} \big| f^0\big(X_n(0;t+h,x+k,v),\left\lvert V_n(0;t+h,x+k,v)\right\rvert\big) \\
& \hspace*{2.5in}-f^0\big(X_n(0;t,x,v),\left\lvert V_n(0;t,x,v)\right\rvert\big) \big| dv < \frac{\eta }{4}, \\
&(\text{resp. } \\
&\sup_{(t,\,x)\,\in\,[0,\,T]\,\times\,[0,\,1]} \int _{|v|\,<\,U} \big| g^0\big(Y_n(0;t+h,x+k,v),\left\lvert W_n(0;t+h,x+k,v)\right\rvert\big) \\
&\hspace*{2.5in}-g^0\big(Y_n(0;t,x,v),\lvert W_n(0;t,x,v)\rvert\big) \big| dv < \frac{\eta }{4}, \\
&\hspace*{3.83in}n \in \N,\quad |h|+|k|<h_0,
\end{align*}
follows from the uniform continuity on $[0,T] \times [0,1] \times [-U,U]$ of the map
\[
(t,x,v) \mapsto \big (X_n(0;t,x,v), \lvert V_n(0;t,x,v)\rvert, Y_n(0;t,x,v), \lvert W_n(0;t,x,v) \rvert \big),
\]
and its continuous dependence with respect to the fields. The Ascoli theorem applies, which ends the proof of the second step.
\end{proof}


\begin{proof}[{Third step of the proof}]
Let us prove that $S$ is continuous in $C([0,T] \times [0,1])$. Let $(a_n)_{n\in \N }$ converging to $a$ in $C([0,T] \times [0,1])$. Denote by
\begin{align*}
E_n &= S_1(a_n), \quad (f_n,g_n) = (S_2 \circ S_1) (a_n), \quad S(a_n) = \int_{\R } (f_n-g_n)dv,\\
E &= S_1(a),\mkern41mu (f,g) = (S_2 \circ S_1)(a), \mkern35mu S(a) = \int_{\R } (f-g)dv.
\end{align*}
The sequence $(E_n)_{n\,\in\,\N}$ converges to $E$ in $C([0,T] \times [0,1])$, because
\[
\max_{[0,\,T]\,\times\,[0,\,1]} \left|E_n-E\right| \leq \max_{[0,\,T]\,\times\,[0,\,1]} \left|a_n-a\right|.
\]
Let $\eta>0$ be given. By~\eqref{hyp1-f0}, there is $U>0$ such that,
\begin{align*}
&\int_{|v|\,>\,U} \sup_{x\,\in\,[0,\,1],\,|w-v|\,<\left(\left|E_0\right|+M\right)T} f^0(x,w)dv \\
+ &\int_{|v|\,>\,U} \sup_{x\,\in\,[0,\,1],\,|w-v|\,<\,\left(\left|E_0\right|+M\right)T} g^0(x,w)dv < \eta.
\end{align*}
And so,
\begin{multline*}
\int_{|v|\,>\,U} \Big(f^0\big(X_n(0;t,x,v),\lvert V_n(0;t,x,v)\rvert\big)+g^0\big(Y_n(0;t,x,v),\lvert W_n(0;t,x,v)\rvert\big)\Big) dv\\
+\int_{|v|\,>\,U} \Big(f^0\big(X(0;t,x,v),\lvert V(0;t,x,v)\rvert\big)+g^0\big(Y(0;t,x,v),\lvert W(0;t,x,v)\rvert\big)\big) dv\\
\leq 4\eta,\quad n\in \N ^*.
\end{multline*}
The convergence of
\begin{align*}
&\int_{|v|\,<\,U} \Big(f^0\big(X_n(0;t,x,v),|V_n(0;t,x,v)|\big)-g^0\big(Y_n(0;t,x,v),|W_n(0;t,x,v)|\big)\Big) dv
\\
\intertext{to}
&\int_{|v|\,<\,U} \Big(f^0\big(X(0;t,x,v),|V(0;t,x,v)|\big)-g^0\big(Y(0;t,x,v),|W(0;t,x,v)|\big)\Big) dv
\end{align*}
in $C([0,T] \times [0,1])$ when $n \to +\infty $ follows from the continuous dependence of
\begin{multline*}
[0,T] \times [0,1] \times [-U,U] \ni (t,x,v)\\ \mapsto \big(X(0;t,x,v),|V(0;t,x,v)|,Y(0;t,x,v),|W(0;t,x,v)|\big)
\end{multline*}
with respect to the fields.
\end{proof}



\section{Proof of the uniqueness Theorem~\ref{main-result-unicite}}\label{pf-main-result-unicite}

This section splits into two lemmas. Under the assumptions of Theorem~\ref{main-result-unicite} and locally in time, Lemma~\ref{cadre-impliquant-nb-fini-de-rebonds} provides a uniform bound with respect to $(t,x,v)\in [0,T] \times [0,1] \times \R $ of the number of possible bounces at the boundary of the domain of the generalized characteristics associated to a solution to~\eqref{VP-f}-\eqref{boundaries-E}. Lemma~\ref{dm-unicite} proves the local in time uniqueness result of Theorem~\ref{main-result-unicite}.

\begin{lemm}\label{cadre-impliquant-nb-fini-de-rebonds}
Assume that for some $c_0>0$ and $V_0>c_0$,
\begin{align}
f^0(x,v) &= 0,\quad x\in [0,1],\quad |v| \leq c_0 \quad \text{or} \quad |v| \geq V_0.\label{troncature2-f0}
\\
\intertext{Let $T\in ] 0,\tfrac{c_0}{2c_i}[$, where}
c_i &= |E_0| + \frac{1}{\epsilon } \int_0^1 \int \left(f^0+g^0\right)(x,v)dvdx.\label{restriction-on-T}
\end{align}
Then the number of possible bounces of the generalized characteristics of a solution to~\eqref{VP-f}--\eqref{boundaries-E} at the boundaries $x=0$ and $x=1$ of the domain is bounded by
\[
\frac{Tc_i}{\sqrt{V_0^2+2c_i}-2V_0} \,.
\]
\end{lemm}


\begin{proof}[{Proof of Lemma~\ref{cadre-impliquant-nb-fini-de-rebonds}}]
By the mass conservations,
\[
\int f(t,x,v)dxdv= \int f^0(x,v)dxdv, \quad \int g(t,x,v)dxdv= \int g^0(x,v)dxdv, \quad t\in [0,T].
\]
Hence, $c_i$ is a bound from above of $\|E\|_{\infty}$. Let $x_0 \in [0,1]$ and $c_0<|v_0|<V_0$. By~\eqref{property-lemma-V-bounded} and~\eqref{restriction-on-T}, the velocities $V(s;0,x_0,v_0)$ along the generalized characteristics starting at $(0,x_0,v_0)$ satisfy
\begin{equation}\label{pf-lemma4.1-a}
\frac{c_0}{2}\leq |V(s;0,x_0,v_0)| \leq 2V_0,\quad s\in [0,T].
\end{equation}
Assume a first bounce on the generalized characteristics starting at $(0,x_0,v_0)$ occurs at $s_1$ on the $x=1$ boundary. It follows from the positive bound from below of $|V(s;0,x_0,v_0)|$ in~\eqref{pf-lemma4.1-a} that the next possible bounce will occur at $x=0$. Denote by $s_2$ the time of such a bounce. Since
\begin{gather*}
X\left(s_2;0,x_0,v_0\right) = X\left(s_2;s_1,1,V\left(s_1^+;0,x_0,v_0\right)\right),
\\
\intertext{it holds that}
1+V\left(s_1^+;0,x_0,v_0\right)\left(s_2-s_1\right)+\int_{s_1}^{s_2}\left(s_2-r\right)E\left(r,X\left(r;0,x_0,v_0\right)\right)dr = 0.
\end{gather*}
Taking into account that $| V(s_1^+;0,x_0,v_0)| $ (resp. $E(r,X(r;0,x_0,v_0))$) is bounded from above (resp. from below) by $2V_0$ (resp. $-c_i$) implies that
\[
c_i\left(s_2-s_1\right)^2+2V_0\left(s_2-s_1\right)-2\geq 0.
\]
The result of the lemma follows.
\end{proof}

\begin{lemm}\label{dm-unicite}
Under the assumptions of Theorem~\ref{main-result-unicite}, the solution to the Cauchy problem~\eqref{VP-f}-\eqref{boundaries-E} is locally unique in time.
\end{lemm}
\begin{proof}[{Proof of Lemma~\ref{dm-unicite}}]
Let $T>0$ be given. Let $(f,g)$ and $(\tilde{f},\tilde{g})$ be two solutions to the problem on $[0,T]$. Denote by $(X(\cdot ; t,x,v),V(\cdot ; t,x,v))$ (resp.
\begin{align*}
\big(Y(\cdot ; t,x,v),W(\cdot ; t,x,v)\big), \quad &\text{resp.} \left(\tilde{X}(\cdot ; t,x,v),\tilde{V}(\cdot ; t,x,v)\right), \\ 
&\text{resp.} \left(\tilde{Y}(\cdot ; t,x,v),\tilde{W}(\cdot ; t,x,v)\right),
\end{align*}
the generalized characteristics associated to $f$ (resp. $g$, resp. $\tilde{f}$, resp. $\tilde{g}$) starting at $(t,x,v)$. Let us prove that for $T_0\in]0,T[$ small enough,
\begin{multline*}
\big(X(0;t,x,v),|V(0;t,x,v)|,Y(0;t,x,v),|W(0;t,x,v)|\big)\\
= \big(\tilde{X}(0;t,x,v),\left|\tilde{V}(0;t,x,v)\right|,\tilde{Y}(0;t,x,v),\left|\tilde{W}(0;t,x,v)\right|\big),\\
(t,x,v)\in [0,T_0] \times [0,1] \times \R.
\end{multline*}
By Lemma~\ref{cadre-impliquant-nb-fini-de-rebonds}, the number of bounces on $[0,T]$ of generalized characteristics on the boundaries $x=0$ and $x=1$ is uniformly bounded. Let us first consider the case where the generalized characteristics encounter no bounce (resp. at most one bounce) at $x=0$ (resp. at $x=1$) on $[0,T]$. Denote by $A_0$ (resp. $A_2$) the set of $(r,y,u)\in [0,T] \times [0,1] \times \R $ such that the $(X,V)$ and $(\tilde{X},\tilde{V})$ characteristics respectively associated to $f$ and $\tilde{f}$ and passing at $(r,y,u)$, both have no bounce (resp. one bounce) at $x= 1$ on $[0,T]$. Denote by $A_1$ the set of $(r,y,u)\in [0,T] \times [0,1] \times \R $ such that the $(X,V)$ (resp. the $(\tilde{X},\tilde{V})$) characteristics passing at $(r,y,u)$ has one bounce (resp. no bounce) at $x= 1$ on $[0,T]$. Denote by $\bar{A}_1$ the set of $(r,y,u)\in [0,T] \times [0,1] \times \R $ such that the $(X,V)$ (resp. the $(\tilde{X},\tilde{V})$) characteristics passing at $(r,y,u)$ has no bounce (resp. one bounce) at $x= 1$ on $[0,T]$. Denote by $(B_i)_{0\leq i\leq 2}\cup \bar{B}_1$ analogous sets relative to the $(Y,W)$ and $(\tilde{Y},\tilde{W})$ characteristics respectively associated to $g$ and $\tilde{g}$. Denote by
\begin{align*}
\alpha &= \sup _{\left(r,\,y,\,u\right)\,\in\, A_0\,\cup\,A_2}\Big(\left| \left(X-\tilde{X}\right)(0;r,y,u)\right| +\left| \left(V-\tilde{V}\right)(0;r,y,u)\right| \Big) \\
&+\sup _{\left(r,\,y,\,u\right)\,\in\,A_1\,\cup \,\bar{A}_1}\Big(\left| \left(X-\tilde{X}\right)(0;r,y,u)\right| +\left| \left(V+\tilde{V}\right)(0;r,y,u)\right| \Big) \\
&+\sup _{\left(r,\,y,\,u\right)\,\in\,B_0\,\cup \,B_2}\Big(\left| \left(Y-\tilde{Y}\right)(0;r,y,u)\right| +\left| \left(W-\tilde{W}\right)(0;r,y,u)\right| \Big) \\
&+\sup _{\left(r,\,y,\,u\right)\,\in\,B_1}\Big(\left| \left(Y-\tilde{Y}\right)(0;r,y,u)\right| +\left| \left(W+\tilde{W}\right)(0;r,y,u)\right| \Big).
\end{align*}
Using the Lipschitz assumption on $f_0$ and $g_0$, and their evenness when considering $A_1$ and $B_1$, notice that
\begin{multline*}
\left|E-\tilde{E}\right|(r,z)\\
\begin{aligned}
&= \Bigg | \int_0^z\int \left[f^0\big(X(0; r,y,u),V(0;r,y,u)\big)-f ^0\left(\tilde{X}(0;r,y,u),\tilde{V}(0;r,y,u)\right) \right] dudy \\
&-\int_0^z\int \left [g^0\big(Y(0; r,y,u),W(0;r,y,u)\big)-g ^0\left(\tilde{Y}(0;r,y,u),\tilde{W}(0;r,y,u)\right) \right] dudy \Bigg | \\
&\leq c\alpha, \quad (r,z) \in [0,T] \times [0,1].
\end{aligned}
\end{multline*}
Let $(t,x,v)\in A_0$. Both generalized characteristics $(X,V)$ and $(\tilde{X},\tilde{V})$ without backward bounce, starting at $(t,x,v)$, are given by
\begin{align*}
&X(s;t,x,v)= x+v(s-t)+\int_s^t \big(r-s)E(r,X(r;t,x,v)\big)dr, \\
&V(s;t,x,v)= v-\int_s^t E\big(r,X(r;t,x,v)\big)dr, \\
&\tilde{X}(s;t,x,v)= x+v(s-t)+\int_s^t (r-s)\tilde{E}\left(r,\tilde{X}(r;t,x,v)\right)dr, \\
&\tilde{V}(s;t,x,v)= v-\int_s^t \tilde{E}\left(r,\tilde{X}(r;t,x,v)\right)dr,\quad s\in [0,t].
\end{align*}
Hence

\begin{multline*}
\left (\left|X-\tilde{X}\right| + \left|V-\tilde{V}\right| \right) (s;t,x,v) \\
\begin{aligned}
&\leq c \left(\int_s^t \left|E-\tilde{E}\right|\left(r,\tilde{X}(r;t,x,v)\right)dr\right.\\
&\quad\; \left.+ \int_s^t \left | E(r,X(r;t,x,v)-E\left(r,\tilde{X}(r;t,x,v)\right) \right| dr \right) \\
&\leq c \left(\int_s^t \left|E-\tilde{E}\right| \left(r,\tilde{X}(r;t,x,v)\right) dr + \int_s^t \left|X-\tilde{X}\right| (r;t,x,v) dr \right) \\
&\leq c \left(T\alpha + \int_s^t \left|X-\tilde{X}\right| (r;t,x,v) dr \right),\quad s\in [0,t].
\end{aligned}
\end{multline*}
Consequently,
\[
\left (\left|X-\tilde{X}\right| +\left|V-\tilde{V}\right| \right) (s;t,x,v) \leq cT\alpha, \quad s \in [0,T], \quad (t,x,v)\in A_0.
\]
In particular,
\begin{equation}\label{uniq-3}
\left (\left|X-\tilde{X}\right|+\left|V-\tilde{V}\right| \right)(0;t,x,v) \leq cT\alpha, \quad (t,x,v) \in A_0.
\end{equation}


Let $(t,x,v)\in A_2$. Both generalized characteristics with a backward bounce at $(t_1(t,x,v),1)$ for $(X,V)$ (resp. at $(\tilde{t}_1(t,x,v),1)$ for $(\tilde{X},\tilde{V})$), starting at $(t,x,v)$, are given by
\begin{align*}
X(s;t,x,v)&= x+v(s-t)+\int_s^t (r-s)E\left(r,X(r;t,x,v)\right)dr, \\
V(s;t,x,v)&= v-\int_s^t E(r,X(r;t,x,v))dr, \quad s \in [t_1(t,x,v),t], \\
X(s;t,x,v)&= 1+(t_1-s) \left (v-\int_{t_1}^t E(r,X(r;t,x,v))dr \right) \\
&\qquad + \int_s^{t_1} (r-s)E(r,X(r;t,x,v))dr, \\
V(s;t,x,v) &= -v+\int_{t_1}^t E(r,X(r;t,x,v))dr 
\\
&\qquad- \int_s^{t_1} E(r,X(r;t,x,v))dr, \; s \in [0,t_1(t,x,v)[,
\end{align*}
and similar equations for $(\tilde{X},\tilde{V})$. Assume $t_1(t,x,v)\leq
\tilde{t}_1(t,x,v)$. It holds that
\begin{align*}
\left(X-\tilde{X}\right)(s) &= \int_s^t (r-s) \left [E(r,X(r))-\tilde{E}\left(r,\tilde{X}(r)\right) \right] dr, \\
\left(V-\tilde{V}\right)(s) &= \int_s^t \left[E(r,X(r))-\tilde{E}\left(r,\tilde{X}(r)\right) \right] dr, \quad s \in \left[\tilde{t}_1(t,x,v),t\right].
\end{align*}
Hence,
\begin{equation}\label{A2-a}
\left (\left|X-\tilde{X}\right|+\left|V-\tilde{V}\right| \right) (s;t,x,v) \leq cT\alpha, \quad s \in \left[\tilde{t}_1(t,x,v),t\right], \quad (t,x,v)\in A_2,
\end{equation}
as in the $A_0$ case. Moreover,
\begin{multline}\label{A2-b}
\left(X-\tilde{X}\right)(s)\\
\begin{aligned}
&= 2\left(\tilde{t}_1-s\right) \left(-v + \int _{\tilde{t}_1}^t \tilde{E}\left(r,\tilde{X}(r)\right)dr \right) + \int_s^{\tilde{t}_1} (r-s) (E(r,X(r))-\tilde{E}\left(r,\tilde{X}(r)\right) dr \\
&+ \int_{\tilde{t}_1}^t (r-s) (E(r,X(r))-\tilde{E}\left(r,\tilde{X}(r)\right)dr, \qquad s \in \left[t_1(t,x,v), \tilde{t}_1(t,x,v)\right[.
\end{aligned}
\end{multline}
The distance from $t_1$ to $\tilde{t}_1$ can be controlled in the following way. The definition of $(t_1(t,x,v), V(t_1^+))$,
\[
x-1+v(t_1-t)+\int_{t_1}^t (r-t_1)E(r,X(r))dr = 0, \quad v-\int_{t_1}^t E(r,X(r))dr = V(t_1^+),
\]
implies that
\begin{multline*}
V\left(t_1^-\right)\left(\tilde{t}_1-t_1\right)\\
= \int_{\tilde{t}_1}^{t_1} \left(r-\tilde{t}_1\right)E(r,X(r))dr + \int_{\tilde{t}_1}^t \left(r-\tilde{t}_1\right) \left(\tilde{E}\left(r,\tilde{X}(r)\right)-E(r,X(r)) \right) dr.
\end{multline*}
Hence, for $T$ small enough,
\begin{equation}\label{t1-tTilde1}
\begin{split}
0\leq \tilde{t}_1-t_1&\leq cT\left(\sup_{r\,\in\, \left[\tilde{t}_1,\,t\right]} \left|X-\tilde{X}\right|(r)+\alpha \right) \\
&\leq cT\alpha,
\end{split}
\end{equation}
by~\eqref{A2-a}. Consequently,
\begin{equation}\label{A2-c}
\left|X-\tilde{X}\right|(s) \leq c\left|\tilde{t}_1-t_1\right| \leq cT\alpha, \quad s \in \left[t_1(t,x,v),\tilde{t}_1(t,x,v)\right[.
\end{equation}
Finally,
\begin{multline*}
\left(X-\tilde{X}\right)(s)\\
\begin{aligned}
&= v\left(t_1-\tilde{t}_1\right)-(t_1-s)\int_{t_1}^t E(r,X(r))dr + \left(\tilde{t}_1-s\right)\int_{\tilde{t}_1}^t \tilde{E}\left(r,\tilde{X}(r)\right)dr \\
&+\int_s^{t_1} (r-s) E(r,X(r))dr - \int_s^{\tilde{t}_1} (r-s)\tilde{E}\left(r,\tilde{X}(r)\right)dr\\
&= \left(\tilde{t}_1-t_1\right) \left(-v + \int_{t_1}^t E(r,X(r))dr \right)\\
&\quad-\left(\tilde{t}_1-s\right) \int_{\tilde{t}_1}^t \left (E(r,X(r))-\tilde{E}\left(r,\tilde{X}(r)\right) \right) dr \\
&- \left(\tilde{t}_1-s\right) \int_{t_1}^{\tilde{t}_1} E(r,X(r)) dr + \int_s^{\tilde{t}_1} (r-s) \left(E(r,X(r))-\tilde{E}\left(r,\tilde{X}(r)\right) \right) dr \\
&+ \int_{\tilde{t}_1}^{t_1} (r-s)E(r,X(r)) dr,
\end{aligned}
\end{multline*}
\begin{align*}
\left(V-\tilde{V}\right)(s) &= 2\int_{t_1}^{\tilde{t}_1} E(r,X(r)) dr + \int_{\tilde{t}_1}^t (E(r,X(r))-\tilde{E}\left(r,\tilde{X}(r) \right) dr \\
&- \int_s^{\tilde{t}_1} \left (E(r,X(r))-\tilde{E}(r,\tilde{X}(r) \right) dr, \qquad s \in [0,t_1[.
\end{align*}
Hence,
\begin{multline}\label{A2-d}
\left|X-\tilde{X}\right|(s) \\
\begin{aligned}
&\leq c\left|\tilde{t}_1-t_1\right|+cT\alpha +cT\sup_{r\,\in\, \left[t_1,\,t\right]} \left|X-\tilde{X}\right|(r) + c\int_s^{t_1} \left|X-\tilde{X}\right|(r) dr \\
&\leq cT\alpha + c\int_s^{t_1} \left|X-\tilde{X}\right|(r) dr, \quad \text{by~\eqref{t1-tTilde1},~\eqref{A2-a} and~\eqref{A2-c}}, \\
&\leq cT\alpha, \quad s \in [0,t_1[, 
\end{aligned}
\end{multline}
\begin{equation}\label{A2-e}
\begin{split}
\left|V-\tilde{V}\right|(s) &\leq c\left|\tilde{t}_1-t_1\right|+cT\alpha +cT\sup_{r\,\in\,[0,\,t]} \left|X-\tilde{X}\right|(r) \\
&\leq cT\alpha, \quad s \in [0,t_1[, \quad \text{by~\eqref{A2-a}, \eqref{A2-c} and~\eqref{A2-d}}. 
\end{split}
\end{equation}
It follows from~\eqref{A2-d}--\eqref{A2-e} taken at $s=0$ that
\begin{equation}\label{A2-f}
\left (\left|X-\tilde{X}\right|+\left|V-\tilde{V}\right| \right) (0;t,x,v) \leq cT\alpha, \quad (t,x,v) \in A_2.
\end{equation}

Let us consider the last case where $(t,x,v)\in A_1$, the case where $(t,x,v)\in \bar{A}_1$ being similarly treated. Again,
\begin{equation}\label{uniq-5}
\left (\left|X-\tilde{X}\right|+\left|V-\tilde{V}\right| \right) (s;t,x,v) \leq cT\alpha, \quad s \in \left[t_1(t,x,v),t\right].
\end{equation}
Moreover,
\begin{align*}
&X(s)= X\left(t_1\right)+ \left(t_1-s\right)V\left(t_1^+\right)+\int_s^{t_1} (r-s)E(r,X(r)) dr, \\
& V(s)= -V\left(t_1^+\right)-\int_s^{t_1} E(r,X(r)) dr, \\
&\tilde{X}(s)= \tilde{X}\left(t_1\right)-\left(t_1-s\right)\tilde{V}\left(t_1\right)+\int_s^{t_1} (r-s)\tilde{E}\left(r,\tilde{X}(r)\right) dr, \\
&\tilde{V}(s)= \tilde{V}\left(t_1\right)-\int_s^{t_1} \tilde{E}\left(r,\tilde{X}(r)\right) dr,\qquad s\in \left[0,t_1(t,x,v)\right[.
\end{align*}
Consequently,
\begin{multline}\label{uniq-6}
\left (\left|X-\tilde{X}\right|+\left|V+\tilde{V}\right| \right)(0;t,x,v)\\
\begin{aligned}
&\leq \left(\left|X-\tilde{X}\right|+\left|V-\tilde{V}\right| \right)\left(t_1^+(t,x,v); t,x,v\right)+ct_1\\
&\leq cT\alpha +ct_1,\; \text{by }\eqref{uniq-5}.
\end{aligned}
\end{multline}
Moreover, $\tilde{X}$ being non-increasing on $[0,t]$ because the sign of $\tilde{V}$ cannot change,
\[
0 \leq \tilde{X}(0)-\tilde{X}\left(t_1\right)\leq 1-\tilde{X}\left(t_1\right)=\left(X-\tilde{X}\right)\left(t_1\right) \leq cT\alpha.
\]
It then follows from $\tilde{X}(0)-\tilde{X}(t_1)= t_1 |\tilde{V}(\tau)|$ for some $\tau \in [0,t_1]$ that $t_1\leq cT\alpha $. Together with~\eqref{uniq-6}, this implies that
\begin{equation}\label{uniq-7}
\left (\left|X-\tilde{X}\right|+\left|V+\tilde{V}\right| \right)(0;t,x,v) \leq cT\alpha, \quad (t,x,v)\in A_1.
\end{equation}
It follows from~\eqref{uniq-3}, \eqref{A2-f} and~\eqref{uniq-7} that
\begin{align*}
&\sup_{(t,\,x,\,v)\,\in\,A_0\,\cup\,A_2} \left(\left|X-\tilde{X}\right|+\left|V-\tilde{V}\right|\right)(0;t,x,v) \\
+ &\sup_{(t,\,x,\,v)\,\in\,A_1\,\cup\,\bar{A}_1} \left(\left|X-\tilde{X}\right|+\left|V+\tilde{V}\right|\right)(0;t,x,v) \leq cT\alpha.
\end{align*}
It similarly holds that
\begin{align*}
&\sup_{(t,\,x,\,v)\,\in\,B_0\,\cup\,B_2} \left(\left|Y-\tilde{Y}\right|+\left|W-\tilde{W}\right|\right)(0;t,x,v)\\
 + &\sup_{(t,\,x,\,v)\,\in\,B_1\,\cup\,\bar{B}_1} \left(\left|Y-\tilde{Y}\right|+\left|W+\tilde{W}\right|\right)(0;t,x,v) \leq cT\alpha.
\end{align*}
The case where more than zero (resp. one) bounce occurs at $x=0$ (resp. at $x=1$) can be analogously treated by splitting the characteristics between those bouncing an even (resp. odd) number of times. And so, $\alpha \leq cT\alpha $. Hence $\alpha = 0$ for $T$ small enough. Consequently,
\begin{align*}
&(X,Y)(0;t,x,v)= \left(\tilde{X},\tilde{Y}\right)(0;t,x,v),\\
&(V,W)(0;t,x,v)= \pm \left(\tilde{V},\tilde{W}\right)(0;t,x,v),\quad (t,x) \in [0,T] \times [0,1], \quad \frac{c_0}{2} \leq |v| \leq 2V_0.
\end{align*}
It follows that $(f,g)= (\tilde{f},\tilde{g})$ on $[0,T]$.
\end{proof}



\section{Electroneutrality}\label{th-electroneutralite}

The obtention of quasineutrality from the Vlasov--Poisson system, i.e., the passage to the limit when $\epsilon \to 0$ in~\eqref{VP-f}--\eqref{boundaries-E} is a difficult problem. The formal limit does not hold for unstable profiles, as proven by D.~Han-Kwan and M.~Hauray in~\cite{Han15}. In~\cite{Han16}, D.~Han-Kwan and F.~Rousset justified the quasineutral limit of a Vlasov--Poisson system with adiabatic electrons for small times in Sobolev spaces, and for initial data satisfying a Penrose stability condition.


In the following theorem, we pass to the limit when $\epsilon \rightarrow 0$ in~\eqref{E} only, and prove that the electroneutrality equation~\eqref{eq-quasi-neutrality} holds at the limit.

\begin{theo}\label{electroneutralite}
For every $\epsilon >0$ let $(f_{\epsilon },g_{\epsilon },E_{\epsilon })$ be a solution on $[0,T]$ to the Vlasov--Poisson system~\eqref{VP-f}--\eqref{boundaries-E} with initial datum $(f^0_{\epsilon },g^0_{\epsilon },E_{\epsilon,0})$ satisfying the assumptions of Theorem~\ref{main-result-existence}, and
\begin{multline}\label{initial-finite-kinetic-energy}
\int_{[0,\,1]\,\times\,\R } v^2 \left(f_{\epsilon }^0+g_{\epsilon }^0\right) dxdv \leq c, \\
\epsilon >0, \quad \text{(uniformly bounded initial kinetic energy)},
\end{multline}
\begin{equation}\label{uniform-bound-E-boundary}
\sqrt{\epsilon } E_{\epsilon,\,0} \leq c, \quad \epsilon >0, 
\end{equation}
\begin{multline}\label{assumption-initial-electroneutrality}
\int_{\R } f_{\epsilon }^0(x,v) dv
= \int_{\R } g_{\epsilon }^0(x,v) dv, \\ \text{a.e. }x\in [0,1],
\quad \epsilon >0, \;
 \text{(initial electroneutrality)},\quad
\left \| f_{\epsilon }^0 \right \|_{L^2} + \left \| g_{\epsilon }^0 \right \|_{L^2} \leq c, \quad \epsilon >0,
\end{multline}
for some $c>0$. There exist a subsequence $(f_{\epsilon _n}, g_{\epsilon _n})$ of $(f_{\epsilon },g_{\epsilon })$, a subsequence $(E_{\epsilon _n})$ of $(E_{\epsilon })$, and functions $(f,g)$ such that
\begin{align*}
\left(f_{\epsilon _n}, g_{\epsilon _n}\right)_{n\,\in\,\N}\quad &\text{weakly converges to } (f,g) \text{ in } L^2([0,T]\times [0,1] \times \R)\text{ when }n\rightarrow +\infty,\\
\left(\sqrt{\epsilon_n} E_{\epsilon_n}\right)_{n\,\in\, \N} \quad &\text{weakly converges } \text{in}\;L^2([0,T] \times [0,1])\text{ when }n\rightarrow +\infty.
\end{align*}
Moreover, $(\int_{\R } f_{\epsilon _n}dv)_{n\,\in\,\N }$ (resp. $(\int_{\R } g_{\epsilon _n}dv)_{n\,\in\,\N }$) weakly converges in $L^2([0,T] \times [0,1])$ to $\int_{\R } f dv$ (resp. $\int_{\R } g dv$) when $n\rightarrow +\infty $, and
\begin{equation}\label{eq-quasi-neutrality}
\int_{\R } f(t,x,v) dv = \int_{\R } g(t,x,v) dv, \quad \text{a.a. }(t,x) \in [0,T] \times [0,1].
\end{equation}
\end{theo}

\begin{proof}[{Proof of Theorem~\ref{electroneutralite}}]
The proof is classical and relies on the conservation of energy. The energy associated to the Vlasov--Poisson system~\eqref{VP-f}--\eqref{boundaries-E} is the sum of the kinetic and potential energies,
\begin{equation}\label{energy}
\clE_{\epsilon }(t) = \frac{1}{2} \int_{[0,\,1]\, \times\R } v^2 \left(f_{\epsilon } + g_{\epsilon }\right)(t,x,v) dxdv + \frac{\epsilon }{2} \int_{[0,\,1]} E_{\epsilon}(t,x)^2 dx.
\end{equation}
The boundary conditions have been chosen in order to ensure the conservation of the energy. This classically follows from the multiplication by $v^2$ of the Vlasov equations~\eqref{VP-f}--\eqref{VP-g}, their integration with respect to $(x,v)$, and the use of the continuity equations,
\[
\partial_t \left(\int_{\R } \left(f_{\epsilon } -g_{\epsilon }\right) dv \right) + \partial_x \left(\int_{\R } v \left(f_{\epsilon }-g_{\epsilon }\right) dv \right) = 0.
\]
More details can be found in~\cite{Gio19}. And so,
\[
\clE_{\epsilon }(t) = \frac{1}{2}\int_{[0,\,1] \times \R } v^2 \left(f_{\epsilon }^0 + g_{\epsilon }^0\right)(x,v)dxdv + \frac{\epsilon }{2} E_{\epsilon,\,0}^2,\quad t\in [0,T].
\]
The family $(\sqrt{\epsilon } E_{\epsilon })_{\epsilon\,>\,0}$ being uniformly bounded in $L^2([0,T] \times [0,1])$, there is a sequence $(\epsilon _n)$ tending to zero when $n \to +\infty $ such that $(\sqrt{\epsilon _n} E_{\epsilon _n})$ weakly converges in $L^2([0,T] \times [0,1])$. Hence $(\epsilon _n E_{\epsilon _n})$ weakly converges to zero in $L^2([0,T]\times [0,1])$ when $n \to +\infty $. The family $(f_{\epsilon })$ and $(g_{\epsilon })$ being uniformly bounded in $L^2([0,T] \times [0,1] \times\R)$, there is a subsequence of $(\epsilon _n)_{n\,\in\,\N}$, still denoted by $(\epsilon _n)_{n\,\in\,\N}$ for the sake of simplicity, and functions $f$ and $g$ in $L^2([0,T] \times [0,1] \times \R)$, such that $(f_{\epsilon_n})_{n\,\in\,\N}$ (resp. $(g_{\epsilon _n})_{n\,\in\,\N}$) weakly converges in $L^2([0,T] \times [0,1] \times \R)$ to $f$ (resp. $g$) when $n\rightarrow +\infty $, and
\[
\int v^2(f+g)(t,x,v) dxdv < +\infty.
\]
Moreover, $(\int f_{\epsilon _n}(t,x,v) dv)_{n\, \in\,\N} $ (resp. $(\int g_{\epsilon_n}(t,x,v) dv)_{n\,\in\,\N} $) weakly converges in\linebreak $L^2([0,T]\times [0,1])$ to $\int f(t,x,v) dv$ (resp. $\int g(t,x,v) dv$). Indeed, for any function $\alpha \in L^2([0,T] \times [0,1])$,
\begin{multline}\label{weak-cv-of-rho-n}
\left |\int \alpha(t,x) \int \left(f_{\epsilon _n}-f\right)(t,x,v) dvdxdt \right|\\
\leq \frac{\|\alpha\|_{\infty}}{K^2} \int v^2\left(f_{\epsilon _n}+f\right)(t,x,v) dvdxdt + \left| \int \alpha(t,x) \chi_K(v)\left(f_{\epsilon _n}-f\right)(t,x,v) dvdxdt \right|,
\end{multline}
where $\chi_K$ denotes the characteristic function of $]-K,K[$. The first term in the r.h.s. of~\eqref{weak-cv-of-rho-n} tends to zero when $K \to +\infty $ uniformly with respect to $n$. The second term in the r.h.s. of~\eqref{weak-cv-of-rho-n} tends to zero for any fixed $K$, given that the map $\alpha(t,x) \chi_K(v)$ belongs to $L^2([0,T] \times [0,1] \times \R)$. The passage to the limit when $n \to +\infty $ in~\eqref{E} leads to the electroneutrality equation~\eqref{eq-quasi-neutrality}.
\end{proof}
\vspace{30pt}


\begin{thebibliography}{99}
\bibitem{Ale93}Alexandre, Radjesvarane Weak solutions of the Vlasov–Poisson initial boundary value problem, Math. Methods Appl. Sci., Volume 16 (1993) no. 8, pp. 533-607 | MR | Zbl
\bibitem{Ars75}Arsenev, Alekseĭ Existence in the large of a weak solution of Vlasov’s system of equations, Zh. Vychisl. Mat. Mat. Fiz., Volume 15 (1975), pp. 136-147 | MR | Zbl
\bibitem{Ben94}Ben Abdallah, Naoufel Weak solutions of the initial-boundary value problem for the Vlasov–Poisson system, Math. Methods Appl. Sci., Volume 17 (1994) no. 6, pp. 451-476 | MR | Zbl
\bibitem{BD85}Bardos, Claude W.; Degond, Pierre Global existence for the Vlasov–Poisson equation in 3 space variables with small initial data, Ann. Inst. Henri Poincaré, Anal. Non Linéaire, Volume 2 (1985) no. 2, pp. 101-118 | Numdam | MR | Zbl
\bibitem{DL88}DiPerna, Ronald J.; Lions, Pierre-Louis Solutions globales du type Vlasov–Poisson, C. R. Acad. Sci., Paris, Sér. I, Volume 307 (1988) no. 12, pp. 655-658 | MR | Zbl
\bibitem{Gio19}Giorgi, Pierre-Antoine Analyse mathématique de modèles cinétiques en physique des plasmas, Ph. D. Thesis, Aix-Marseille University, Marseille, France (2019)
\bibitem{Guo94}Guo, Yan Regularity for the Vlasov Equation with boundary in a Half Space, Indiana Univ. Math. J., Volume 43 (1994) no. 1, pp. 255-319 | Zbl
\bibitem{Guo95}Guo, Yan Singular solutions of the Vlasov–Maxwell system on a half line, Arch. Ration. Mech. Anal., Volume 131 (1995) no. 3, pp. 241-304 | MR | Zbl
\bibitem{Han15}Han-Kwan, Daniel; Hauray, Maxime Stability issues in the quasineutral limit of the one-dimensional Vlasov–Poisson equation, Commun. Math. Phys., Volume 334 (2015) no. 2, pp. 1101-1152 | MR | DOI | Zbl
\bibitem{Han16}Han-Kwan, Daniel; Rousset, Frédéric Quasineutral limit for Vlasov-Poisson with Penrose stable data, Ann. Sci. Éc. Norm. Supér., Volume 49 (2016) no. 6, pp. 1445-1495 | MR | DOI | Zbl
\bibitem{HM18}Holding, Thomas; Miot, Evelyne Uniqueness and stability for the Vlasov–Poisson system with spatial density in Orlicz spaces, Mathematical analysis in fluid mechanics: selected recent results. International conference on vorticity, rotation and symmetry (IV) – complex fluids and the issue of regularity, CIRM, Luminy, Marseille, France, May 8–12, 2017. Proceedings (Danchin, Raphaël et al., eds.) (Contemporary Mathematics), Volume 710, American Mathematical Society, 2018, pp. 145-162 | MR | Zbl | DOI
\bibitem{HS08}Hwang, Hyung J.; Schaeffer, Jack Uniqueness for weak solutions of a one-dimensional boundary value problem for the Vlasov–Poisson system, J. Differ. Equations, Volume 244 (2008) no. 10, pp. 2665-2691 | MR | DOI | Zbl
\bibitem{Loe06}Loeper, Grégoire Uniqueness of the solution to the Vlasov-Poisson system with bounded density, J. Math. Pures Appl., Volume 86 (2006) no. 1, pp. 68-79 | MR | DOI | Zbl
\bibitem{LP91}Lions, Pierre-Louis; Perthame, Benoît Propagation of moments and regularity for the 3-dimensional Vlasov–Poisson system, Invent. Math., Volume 105 (1991) no. 2, pp. 415-430 | MR | DOI | Zbl
\bibitem{Mio16}Miot, Evelyne A uniqueness criterion for unbounded solutions to the Vlasov–Poisson system, Commun. Math. Phys., Volume 346 (2016) no. 2, pp. 469-482 | MR | DOI | Zbl
\bibitem{Mis99}Mischler, Stéphane On the trace problem for solutions of the Vlasov equation, Commun. Partial Differ. Equations, Volume 25 (1999) no. 7-8, pp. 1415-1443 | MR | DOI | Zbl
\bibitem{Pal12}Pallard, Christophe Moment propagation for weak solutions to the Vlasov–Poisson system, Commun. Partial Differ. Equations, Volume 37 (2012) no. 7, pp. 1273-1285 | MR | DOI | Zbl
\bibitem{Pfa92}Pfaffelmoser, Klaus Global classical solutions of the Vlasov–Poisson system in three dimensions for general initial data, J. Differ. Equations, Volume 95 (1992) no. 2, pp. 281-303 | MR | DOI | Zbl
\bibitem{Pou90}Poupaud, Frédéric Stationary solutions of Vlasov–Poisson equations, C. R. Acad. Sci., Paris, Sér. I, Volume 311 (1990) no. 6, pp. 307-312 | MR | Zbl
\bibitem{Rei97}Rein, Gerhard Self-gravitating systems in Newtonian theory – The Vlasov–Poisson system, Mathematics of gravitation, Part I (Warsaw, 1996) (Banach Center Publications), Volume 41, Polish Academy of Sciences, 1997, pp. 179-194 | Zbl | MR
\bibitem{Rob97}Robert, Raoul Unicité de la solution faible à support compact de l’équation de Vlasov–Poisson, C. R. Acad. Sci., Paris, Sér. I, Volume 324 (1997) no. 8, pp. 833-877 | Zbl
\bibitem{Sch91}Schaeffer, Jack Global existence of smooth solutions to the Vlasov–Poisson system in three dimensions, Commun. Partial Differ. Equations, Volume 16 (1991), pp. 1313-1335 | MR | DOI | Zbl
\bibitem{Wec95}Weckler, Jürge On the initial-boundary-value problem for the Vlasov–Poisson system: existence of weak solutions and stability, Arch. Ration. Mech. Anal., Volume 130 (1995) no. 2, pp. 145-161 | MR | DOI | Zbl
\end{thebibliography}
\end{document}
